\section{Logbook Entry 014: Variance Scaling, Lag Ablation, and Descriptive Ripeness in Spontaneous Allen Neural HRSM}
\label{sec:logbook-entry-014-variance-ablation-ripeness}

\subsection{Purpose}

This entry records the pre-manuscript robustness batch for the spontaneous-focused real Allen Neural HRSM branch. After three independent Allen Neuropixels sessions had been processed, the next question was whether the apparent non-Markovian memory signal survived standard sanity checks.

Three checks were performed:

\begin{enumerate}
    \item a variance-scaling audit, asking whether controlled memory gain merely tracks target variance, dispersion, or lag-1 autocorrelation;
    \item a lag-ablation grid, asking whether the result depends narrowly on the chosen lag pair \((1,2)\);
    \item a descriptive spontaneous neural ripeness index, asking whether session-region states can be ranked by combined HRSM state potential and controlled organizational memory.
\end{enumerate}

These analyses used the three spontaneous-focused Allen sessions:

\begin{verbatim}
715093703
719161530
750749662
\end{verbatim}

Each session used the same spontaneous-focused extraction scale: 20 units per region, 15 spontaneous presentations, 0.25 second bins, and four regions: VISp, VISl, LGd, and CA1.

\subsection{Updated Cross-Session Synthesis}

Before the robustness tests, the cross-session synthesis was updated to include all three sessions using:

\begin{verbatim}
scripts/allen/16_cross_session_spontaneous_synthesis.py
\end{verbatim}

The three-session synthesis showed that all sessions retained numerical HRSM axis separation:

\begin{verbatim}
715093703 max_abs_offdiag_corr = 3.994795e-15
719161530 max_abs_offdiag_corr = 4.175386e-15
750749662 max_abs_offdiag_corr = 6.142750e-15
\end{verbatim}

The session-level summary showed a qualified replication. The first two sessions had positive raw mean-rate memory gain in all four regions, while the third session had one slightly negative raw VISp mean-rate gain. However, all three sessions retained positive observed-minus-shuffled controlled memory at the session level:

\begin{center}
\begin{tabular}{lrrrr}
\hline
Session & Median memory gain & Mean memory gain & Fraction positive gain & Median observed minus shuffled \\
\hline
715093703 & 0.035902 & 0.040904 & 1.00 & 0.039595 \\
719161530 & 0.037543 & 0.032870 & 1.00 & 0.041368 \\
750749662 & 0.030390 & 0.027334 & 0.75 & 0.033282 \\
\hline
\end{tabular}
\end{center}

The three-session replication flags were:

\begin{verbatim}
all_sessions_memory_gain_positive = False
all_sessions_above_shuffle = True
orthogonality_numeric_precision = True
top_targets_are_recruitment_and_entropy = True
LGd_best_phi_in_each_session = False
\end{verbatim}

This means the broadest raw-memory and LGd-dominance claims do not survive unchanged. The stronger controlled-memory claim does survive: observed aligned-history models exceed shuffled-lag controls across sessions, and the top memory-sensitive targets remain active-unit fraction and population entropy.

\subsection{Variance-Scaling Audit}

The variance-scaling audit was run using:

\begin{verbatim}
scripts/allen/17_cross_session_variance_scaling_audit.py
\end{verbatim}

The purpose was to test whether the controlled memory signal was simply a consequence of target variance, target dispersion, or lag-1 autocorrelation. The script joined each session's population-state matrix with the target-level memory sweep and computed target mean, variance, standard deviation, interquartile range, coefficient of variation, and lag-1 autocorrelation.

The raw controlled-memory ranking remained:

\begin{center}
\begin{tabular}{lrr}
\hline
Target & Median observed minus shuffled & Raw rank \\
\hline
active\_unit\_fraction & 0.063814 & 1 \\
population\_rate\_entropy & 0.062511 & 2 \\
population\_mean\_rate\_hz & 0.036210 & 3 \\
population\_std\_rate\_hz & 0.033441 & 4 \\
population\_l2\_rate\_norm & 0.033028 & 5 \\
population\_state\_speed & 0.012205 & 6 \\
\hline
\end{tabular}
\end{center}

After residualizing target variance, lag-1 autocorrelation, and row count, the ranking changed:

\begin{center}
\begin{tabular}{lrr}
\hline
Target & Median residual gain & Residual rank \\
\hline
population\_l2\_rate\_norm & 0.011099 & 1 \\
population\_rate\_entropy & 0.007370 & 2 \\
active\_unit\_fraction & 0.001711 & 3 \\
population\_state\_speed & -0.002276 & 4 \\
population\_mean\_rate\_hz & -0.009043 & 5 \\
population\_std\_rate\_hz & -0.009630 & 6 \\
\hline
\end{tabular}
\end{center}

The variance-scaling flags were:

\begin{verbatim}
top_two_raw_are_recruitment_and_entropy = True
top_two_residual_are_recruitment_and_entropy = False
overall_variance_correlation_not_extreme = True
\end{verbatim}

The overall Spearman correlation between target variance and controlled gain was:

\begin{verbatim}
Spearman(target_var, controlled_gain) = -0.5594571998199241
\end{verbatim}

\subsection{Interpretation of Variance Scaling}

The variance-scaling result is a partial constraint, not a collapse. The raw controlled-memory signal remains strongest for active-unit fraction and entropy. However, after residualization, active-unit fraction no longer remains in the top two; population L2 rate norm rises to first, and entropy remains second.

The correct interpretation is therefore narrower: the recruitment-memory effect is partly entangled with variance and autocorrelation structure, while the entropy-memory effect is more robust under this residualization. The manuscript should not claim that active-unit recruitment is fully independent of variance structure. It should instead state that active-unit fraction and entropy are the strongest raw controlled targets, with entropy surviving variance/autocorrelation residualization more cleanly.

\subsection{Lag-Ablation Grid}

The lag-ablation grid was run using:

\begin{verbatim}
scripts/allen/18_spontaneous_memory_lag_ablation_grid.py
\end{verbatim}

The tested lag sets were:

\begin{verbatim}
lag1      = [1]
lag2      = [2]
lag12     = [1, 2]
lag123    = [1, 2, 3]
lag1234   = [1, 2, 3, 4]
\end{verbatim}

The ablation used 50 shuffled controls per group as a screening run.

The cross-session lag-ablation summary showed that active-unit fraction and population entropy remained the top two controlled-memory targets under the canonical lag pair \((1,2)\). More importantly, both remained positive across all tested lag sets:

\begin{center}
\begin{tabular}{lrrrrr}
\hline
Target & lag1 & lag2 & lag12 & lag123 & lag1234 \\
\hline
active\_unit\_fraction & 0.046499 & 0.046029 & 0.066097 & 0.076122 & 0.087106 \\
population\_rate\_entropy & 0.039511 & 0.048006 & 0.064305 & 0.077218 & 0.081646 \\
population\_mean\_rate\_hz & 0.030403 & 0.027874 & 0.039064 & 0.046647 & 0.050860 \\
population\_state\_speed & 0.004257 & 0.012048 & 0.012949 & 0.012580 & 0.012323 \\
\hline
\end{tabular}
\end{center}

The lag-ablation flags were:

\begin{verbatim}
lag12_top_two_are_recruitment_and_entropy = True
recruitment_entropy_positive_all_lagsets = True
\end{verbatim}

\subsection{Interpretation of Lag Ablation}

The lag-ablation result is stronger than the variance-scaling result. The controlled-memory signal does not depend narrowly on choosing lags 1 and 2. Active-unit fraction and population entropy remain positive across lag sets, and their controlled gains increase as additional recent history is added.

This supports the interpretation that the spontaneous neural memory effect is genuinely historical in the operational sense used here: recent past population structure improves prediction relative to shuffled-lag controls.

\subsection{Descriptive Neural Ripeness Index}

A descriptive spontaneous neural ripeness index was then computed using:

\begin{verbatim}
scripts/allen/19_compute_spontaneous_ripeness_index.py
\end{verbatim}

This index was defined as descriptive, not inferential. It combines rank-normalized components:

\begin{enumerate}
    \item \(\Phi_{\rm neural}\),
    \item HRSM stability coordinate \(S\),
    \item controlled mean-rate memory,
    \item active-unit recruitment memory,
    \item entropy memory.
\end{enumerate}

Penalties were applied for negative raw mean-rate memory gain or nonpositive controlled mean-rate memory gain.

The highest-ranked session-region state was:

\begin{verbatim}
715093703 / VISl
ripeness_score = 0.777083
ripeness_class = high
ripeness_mode = memory_carrier_low_phi
\end{verbatim}

The next strongest states were:

\begin{verbatim}
715093703 / CA1        ripeness_score = 0.643750
750749662 / CA1        ripeness_score = 0.633333
719161530 / VISp       ripeness_score = 0.622917
750749662 / VISl       ripeness_score = 0.591667
719161530 / LGd        ripeness_score = 0.554167
\end{verbatim}

The region-level median ripeness ranking was:

\begin{center}
\begin{tabular}{lrrrr}
\hline
Region & Median ripeness & Mean ripeness & Median \(\Phi\) & Region rank \\
\hline
CA1 & 0.633333 & 0.581250 & -0.113984 & 1 \\
VISl & 0.591667 & 0.634028 & -0.067570 & 2 \\
LGd & 0.477083 & 0.485417 & 0.408474 & 3 \\
VISp & 0.381250 & 0.449306 & 0.131802 & 4 \\
\hline
\end{tabular}
\end{center}

The ripeness flags were:

\begin{verbatim}
all_ripeness_scores_finite = True
top_state_has_positive_controlled_memory = True
top_region_not_forced_to_LGd = True
organizational_memory_positive_in_top_half = True
\end{verbatim}

\subsection{Interpretation of Ripeness}

The ripeness index is sane because it does not merely reproduce \(\Phi_{\rm neural}\) or force LGd to the top. Instead, it distinguishes structural state strength from organizational memory carrying capacity.

LGd remains important as a high-\(\Phi\) structural state, but it is not the highest ripeness region once controlled memory and organizational recruitment are included. CA1 and VISl become the top median-ripeness regions because they combine strong controlled organizational memory with moderate or transitional HRSM structure. This reinforces the distinction between state potential and memory-carrying dynamics.

The phrase \emph{ripeness} should therefore be used carefully. It does not mean consciousness, cognition, or readiness for awareness. In this branch it means a session-region state with sufficient HRSM organization and controlled historical predictability to function as a candidate memory-bearing neural state.

\subsection{Overall Conclusion}

The pre-manuscript robustness batch supports a refined conclusion.

The original broad statement would have been too strong: spontaneous neural memory is not uniformly positive in every raw target and region, LGd is not always the highest-\(\Phi\) state, and active-unit recruitment is not fully independent of target variance and autocorrelation.

The stronger statement survives:

\begin{quote}
Across three spontaneous-focused Allen Neuropixels sessions, aligned recent history improves prediction relative to shuffled-lag controls. The most stable raw controlled targets are active-unit fraction and population entropy. Entropy survives variance/autocorrelation residualization more cleanly than active-unit fraction, and both targets remain positive across lag-ablation tests. Descriptive ripeness ranks CA1 and VISl above LGd when controlled organizational memory is included, separating structural HRSM potential from memory-bearing dynamics.
\end{quote}

\subsection{Next Step}

The next step is manuscript extraction. The results section should be built around three claims:

\begin{enumerate}
    \item spontaneous neural HRSM states can be projected reproducibly across Allen sessions with orthogonalized H, R, S, and M axes;
    \item non-Markovian memory is supported relative to shuffled-lag controls, especially for population entropy and active-unit recruitment;
    \item robustness tests narrow the claim: entropy is the most stable organizational-memory target after variance/autocorrelation residualization, while ripeness is descriptive and separates structural potential from memory-carrying dynamics.
\end{enumerate}
